\setcounter{chapter}{13}\chapter{Variables aléatoires, concentration et loi des grands nombres}\label{ch:14}

{\small\sffamily\color{teal}S62, S64, S73, S74, S77\par}

\begin{retenir}{Objectifs}Calculer espérance et variance d'une somme et d'une moyenne ; appliquer Bienaymé-Tchebychev ; déterminer une taille d'échantillon.\end{retenir}

\section{Notes de cours}
\subsection{Sommes et transformations}
Pour une variable aléatoire finie $X$, prenant les valeurs $x_i$ avec probabilités $p_i$,
\[E(X)=\sum_i x_ip_i,\qquad V(X)=\sum_i(x_i-E(X))^2p_i=E(X^2)-E(X)^2.\]
L'écart type est $\sigma(X)=\sqrt{V(X)}$. En développant le carré dans $V(aX+b)$, le décalage $b$ s'élimine et le facteur $a$ est élevé au carré.
Pour des variables aléatoires finies $X,Y$ et des réels $a,b$,
\[E(X+Y)=E(X)+E(Y),\quad E(aX+b)=aE(X)+b,\quad V(aX+b)=a^2V(X).\]
La linéarité de l'espérance ne demande pas l'indépendance. Si $X,Y$ sont indépendantes, $V(X+Y)=V(X)+V(Y)$ (admis). Sans cette hypothèse, on ne peut pas appliquer cette formule. L'écart type de $aX+b$ vaut $|a|\sigma(X)$.
\paragraph{Exemple de dépendance.} Si $Y=X$, alors $V(X+Y)=V(2X)=4V(X)$, et non $2V(X)$ dès que $V(X)\ne0$.
\paragraph{Démonstration pour la binomiale.} Écrire $X=B_1+\cdots+B_n$ avec les indicatrices indépendantes des succès. Chaque $B_i$ vérifie $E(B_i)=p$, $E(B_i^2)=p$ et $V(B_i)=p-p^2$. Par linéarité et additivité sous indépendance, $E(X)=np$, $V(X)=np(1-p)$.
\subsection{Échantillon et moyenne}
Un échantillon de taille $n$ d'une loi est une famille de $n$ variables indépendantes de même loi. Si elles ont espérance $\mu$ et variance $\sigma^2$, on note $S_n=\sum_{i=1}^nX_i$ et $M_n=\frac{S_n}{n}$.
\begin{retenir}{Somme et moyenne}
\[E(S_n)=n\mu,\quad V(S_n)=n\sigma^2,\qquad E(M_n)=\mu,\quad V(M_n)=\frac{\sigma^2}{n},\quad\sigma(M_n)=\frac\sigma{\sqrt n}.\]
\end{retenir}
Pour diviser l'écart type de la moyenne par $2$, multiplier la taille par $4$. La variance de la somme augmente ; celle de la moyenne diminue.
\subsection{Concentration}
\begin{retenir}{Inégalité de Bienaymé-Tchebychev}
Pour une variable d'espérance $\mu$, de variance $\sigma^2$, et $\delta>0$,
\[P(|X-\mu|\ge\delta)\le\frac{\sigma^2}{\delta^2}.\]
Une borne supérieure à $1$ reste vraie, mais n'apporte pas d'information supplémentaire ; on peut la remplacer par $1$.
\end{retenir}
\paragraph{Preuve dans le cadre fini.} Dans $V(X)=\sum_x(x-\mu)^2P(X=x)$, ne conserver que les valeurs vérifiant $|x-\mu|\ge\delta$. Chaque carré est alors au moins $\delta^2$. Donc $V(X)\ge\delta^2P(|X-\mu|\ge\delta)$.
\subsection{Loi des grands nombres}
Appliquée à la moyenne d'un échantillon,
\[P(|M_n-\mu|\ge\delta)\le\frac{\sigma^2}{n\delta^2}\xrightarrow[n\to\infty]{}0.\]
C'est la loi des grands nombres dans le cadre considéré : pour une précision fixée, la probabilité d'un écart au moins égal à cette précision tend vers zéro. Ce résultat ne dit pas que chaque réalisation se rapproche de $\mu$ de façon monotone.
\begin{methode}{Choisir une taille suffisante}
Pour garantir $P(|M_n-\mu|<\delta)\ge1-\alpha$, il suffit de choisir
\[n\ge\frac{\sigma^2}{\alpha\delta^2}.\]
Prendre le plafond. Il s'agit du plus petit entier satisfaisant \emph{cette borne}, pas nécessairement de la taille minimale pour la probabilité exacte.
\end{methode}
Pour une fréquence de succès, $\mu=p$ et $\sigma^2=p(1-p)\le\frac{1}{4}$. Si $p$ est inconnu, $n\ge\frac{1}{4\alpha\delta^2}$ est une garantie uniforme.


\section{Activité de recherche}\textbf{Fluctuations d'une fréquence --- 40 min.}\par On lance indépendamment une pièce équilibrée. Pour tout entier $n\ge1$, on note $F_n$ la fréquence de « pile » parmi les $n$ premiers lancers.
\begin{enumerate}\item Écrire un programme qui simule 1000 lancers, puis représenter les points $(n;F_n)$ pour $1\le n\le1000$. L'écart à $\frac12$ diminue-t-il nécessairement à chaque lancer ?
\item Exprimer $E(F_n)$ et $V(F_n)$ en fonction de $n$.
\item À l'aide de Bienaymé-Tchebychev, majorer $P(|F_{1000}-\frac12|\ge0{,}05)$.
\item Déterminer un entier $n$ suffisant pour que $P(|F_n-\frac12|<0{,}05)\ge0{,}99$.
\item Préciser ce que montre la trajectoire simulée et ce que garantit l'inégalité obtenue. Une seule simulation constitue-t-elle une preuve de cette garantie ?\end{enumerate}

\section{Exercices progressifs}

\exercice{1}{Entraînement}Soit $X$ une variable aléatoire telle que $E(X)=3$ et $V(X)=2$. On pose $Y=4X-1$. Calculer l'espérance, la variance et l'écart type de $Y$.

\exercice{2}{Entraînement}Les variables aléatoires $X$ et $Y$ sont indépendantes et vérifient $E(X)=2$, $E(Y)=5$, $V(X)=1$ et $V(Y)=4$. Calculer $E(2X-Y)$ et $V(2X-Y)$.

\exercice{3}{Entraînement}On considère 64 variables aléatoires indépendantes de même loi, d'espérance $10$ et d'écart type $4$. Calculer l'espérance et l'écart type de leur moyenne.

\exercice{4}{Entraînement}Soit $X$ une variable aléatoire telle que $E(X)=100$ et $V(X)=9$. À l'aide de l'inégalité de Bienaymé-Tchebychev, donner une majoration de $P(|X-100|\ge6)$.

\exercice{5}{Entraînement}On note $M_n$ la moyenne de $n$ variables aléatoires indépendantes de même loi, d'espérance $\mu$ et de variance $4$. Déterminer une taille $n$ suffisante pour que $P(|M_n-\mu|<0{,}2)\ge0{,}95$.

\section{Travail dans le manuel ou une annale}\textbf{Baccalauréat, Métropole, 17 juin 2025 : exercice 1, question 6 (p. 3).} À la maison après S77. Lire le contexte de la question 5 et utiliser les valeurs d'espérance et de variance admises dans la question 6. Le sujet est accessible ici : \href{https://www.education.gouv.fr/sites/default/files/2025-06/baccalaur-at-g-n-ral-2025-math-matiques-440856.pdf}{annale officielle}. Utiliser les données et les consignes de cette annale.

\section{Devoir à la maison}\textbf{DM14 --- Qualité d'une estimation --- 55 min.}\par On répète $n\ge1$ épreuves de Bernoulli indépendantes de même paramètre $p\in[0;1]$, inconnu. On note $F_n$ la fréquence des succès.
\begin{enumerate}\item Démontrer que $E(F_n)=p$ et $V(F_n)=\frac{p(1-p)}n$.
\item Montrer que $p(1-p)\le\frac14$ pour tout $p\in[0;1]$.
\item En déduire une taille $n$ suffisante pour que $P(|F_n-p|<0{,}02)\ge0{,}95$, quelle que soit la valeur de $p$.
\item On sait maintenant que $p\le0{,}1$. Déterminer une majoration plus précise de $p(1-p)$, puis améliorer la taille suffisante obtenue à la question 3.
\item Pour cette question seulement, les $n$ résultats sont tous égaux à une même variable de Bernoulli $B$ de paramètre $p$. Calculer $F_n$ et sa variance. Expliquer pourquoi la garantie précédente ne s'applique plus.\end{enumerate}

\ifprof\clearpage\section{Conduite et corrigés}\begin{profbox}{Repères pour le professeur}Le DS5 et le deuxième blanc précèdent S73 : n'y demander ni Bienaymé-Tchebychev ni loi des grands nombres. Pour la variance, faire entourer l'hypothèse d'indépendance. Pour la concentration, distinguer probabilité exacte et majoration, puis rang suffisant et rang minimal réel. Une simulation n'est pas une preuve et une moyenne observée ne se rapproche pas forcément à chaque nouvel essai.\end{profbox}

\subsection{Corrigé de l'activité}\begin{enumerate}
\item 
\begin{lstlisting}
from random import Random

def trajectoire_piece(N=1000, graine=2026):
    rng = Random(graine)
    piles, frequences = 0, []
    for n in range(1, N+1):
        piles += rng.random() < 0.5
        frequences.append(piles/n)
    return frequences
\end{lstlisting}
Tracer les points $(n,F_n)$ avec la liste obtenue. Après deux essais donnant une pile et une face, $F_2=\frac{1}{2}$. Quel que soit le troisième résultat, $F_3$ vaut $\frac{1}{3}$ ou $\frac{2}{3}$ : l'écart à $\frac{1}{2}$ augmente. La stabilisation n'est pas une diminution monotone de l'écart sur chaque trajectoire.
\begin{center}\begin{tikzpicture}\begin{axis}[width=12cm,height=5cm,xlabel={$n$},ylabel={$F_n$},xmin=1,xmax=1000,ymin=0,ymax=1,grid=major,grid style={gray!15}]
\addplot[teal,thin] coordinates {(1,1.000000) (2,0.500000) (3,0.333333) (4,0.250000) (5,0.400000) (6,0.500000) (7,0.428571) (8,0.375000) (9,0.333333) (10,0.300000) (11,0.272727) (12,0.250000) (13,0.230769) (14,0.214286) (15,0.266667) (16,0.250000) (17,0.294118) (18,0.277778) (19,0.315789) (20,0.300000) (21,0.285714) (22,0.272727) (23,0.304348) (24,0.333333) (25,0.320000) (26,0.307692) (27,0.296296) (28,0.285714) (29,0.275862) (30,0.300000) (31,0.322581) (32,0.312500) (33,0.333333) (34,0.352941) (35,0.342857) (36,0.333333) (37,0.324324) (38,0.342105) (39,0.333333) (40,0.350000) (41,0.365854) (42,0.380952) (43,0.372093) (44,0.386364) (45,0.400000) (46,0.391304) (47,0.382979) (48,0.375000) (49,0.387755) (50,0.380000) (51,0.392157) (52,0.403846) (53,0.396226) (54,0.407407) (55,0.400000) (56,0.410714) (57,0.421053) (58,0.431034) (59,0.440678) (60,0.433333) (61,0.426230) (62,0.435484) (63,0.428571) (64,0.437500) (65,0.446154) (66,0.454545) (67,0.447761) (68,0.441176) (69,0.449275) (70,0.457143) (71,0.450704) (72,0.444444) (73,0.438356) (74,0.432432) (75,0.426667) (76,0.434211) (77,0.428571) (78,0.423077) (79,0.430380) (80,0.425000) (81,0.432099) (82,0.439024) (83,0.445783) (84,0.452381) (85,0.447059) (86,0.441860) (87,0.448276) (88,0.443182) (89,0.449438) (90,0.444444) (91,0.450549) (92,0.456522) (93,0.451613) (94,0.446809) (95,0.442105) (96,0.437500) (97,0.432990) (98,0.428571) (99,0.434343) (100,0.430000) (101,0.435644) (102,0.431373) (103,0.427184) (104,0.423077) (105,0.428571) (106,0.424528) (107,0.420561) (108,0.416667) (109,0.422018) (110,0.427273) (111,0.432432) (112,0.437500) (113,0.442478) (114,0.438596) (115,0.443478) (116,0.448276) (117,0.452991) (118,0.457627) (119,0.453782) (120,0.450000) (121,0.454545) (122,0.450820) (123,0.455285) (124,0.451613) (125,0.448000) (126,0.444444) (127,0.440945) (128,0.437500) (129,0.441860) (130,0.446154) (131,0.442748) (132,0.439394) (133,0.436090) (134,0.440299) (135,0.437037) (136,0.441176) (137,0.445255) (138,0.442029) (139,0.438849) (140,0.435714) (141,0.432624) (142,0.436620) (143,0.433566) (144,0.430556) (145,0.427586) (146,0.424658) (147,0.428571) (148,0.432432) (149,0.429530) (150,0.426667) (151,0.430464) (152,0.434211) (153,0.437908) (154,0.435065) (155,0.438710) (156,0.442308) (157,0.445860) (158,0.443038) (159,0.446541) (160,0.443750) (161,0.447205) (162,0.450617) (163,0.453988) (164,0.457317) (165,0.454545) (166,0.451807) (167,0.449102) (168,0.452381) (169,0.449704) (170,0.452941) (171,0.456140) (172,0.453488) (173,0.450867) (174,0.454023) (175,0.457143) (176,0.460227) (177,0.457627) (178,0.460674) (179,0.458101) (180,0.455556) (181,0.458564) (182,0.456044) (183,0.459016) (184,0.456522) (185,0.454054) (186,0.456989) (187,0.454545) (188,0.452128) (189,0.449735) (190,0.447368) (191,0.450262) (192,0.447917) (193,0.450777) (194,0.453608) (195,0.456410) (196,0.459184) (197,0.461929) (198,0.464646) (199,0.467337) (200,0.470000) (201,0.467662) (202,0.470297) (203,0.472906) (204,0.475490) (205,0.478049) (206,0.480583) (207,0.478261) (208,0.475962) (209,0.473684) (210,0.471429) (211,0.473934) (212,0.471698) (213,0.469484) (214,0.467290) (215,0.469767) (216,0.467593) (217,0.465438) (218,0.463303) (219,0.461187) (220,0.463636) (221,0.466063) (222,0.463964) (223,0.461883) (224,0.459821) (225,0.462222) (226,0.460177) (227,0.458150) (228,0.456140) (229,0.458515) (230,0.460870) (231,0.463203) (232,0.461207) (233,0.463519) (234,0.465812) (235,0.463830) (236,0.466102) (237,0.464135) (238,0.462185) (239,0.460251) (240,0.462500) (241,0.460581) (242,0.458678) (243,0.460905) (244,0.459016) (245,0.461224) (246,0.463415) (247,0.465587) (248,0.467742) (249,0.469880) (250,0.468000) (251,0.466135) (252,0.468254) (253,0.470356) (254,0.472441) (255,0.474510) (256,0.476562) (257,0.474708) (258,0.472868) (259,0.471042) (260,0.473077) (261,0.471264) (262,0.469466) (263,0.467681) (264,0.465909) (265,0.464151) (266,0.466165) (267,0.468165) (268,0.470149) (269,0.468401) (270,0.466667) (271,0.464945) (272,0.463235) (273,0.461538) (274,0.459854) (275,0.461818) (276,0.460145) (277,0.462094) (278,0.460432) (279,0.462366) (280,0.464286) (281,0.462633) (282,0.464539) (283,0.462898) (284,0.461268) (285,0.463158) (286,0.461538) (287,0.463415) (288,0.461806) (289,0.460208) (290,0.462069) (291,0.460481) (292,0.458904) (293,0.460751) (294,0.459184) (295,0.461017) (296,0.462838) (297,0.464646) (298,0.463087) (299,0.461538) (300,0.463333) (301,0.461794) (302,0.463576) (303,0.465347) (304,0.467105) (305,0.468852) (306,0.470588) (307,0.472313) (308,0.470779) (309,0.469256) (310,0.467742) (311,0.466238) (312,0.467949) (313,0.469649) (314,0.471338) (315,0.473016) (316,0.474684) (317,0.476341) (318,0.477987) (319,0.479624) (320,0.478125) (321,0.476636) (322,0.475155) (323,0.473684) (324,0.472222) (325,0.470769) (326,0.469325) (327,0.470948) (328,0.469512) (329,0.468085) (330,0.469697) (331,0.468278) (332,0.469880) (333,0.468468) (334,0.470060) (335,0.468657) (336,0.467262) (337,0.465875) (338,0.464497) (339,0.463127) (340,0.464706) (341,0.463343) (342,0.461988) (343,0.463557) (344,0.462209) (345,0.463768) (346,0.465318) (347,0.466859) (348,0.465517) (349,0.464183) (350,0.465714) (351,0.464387) (352,0.463068) (353,0.464589) (354,0.463277) (355,0.464789) (356,0.463483) (357,0.462185) (358,0.460894) (359,0.459610) (360,0.458333) (361,0.457064) (362,0.455801) (363,0.457300) (364,0.458791) (365,0.457534) (366,0.459016) (367,0.457766) (368,0.456522) (369,0.457995) (370,0.459459) (371,0.458221) (372,0.456989) (373,0.455764) (374,0.454545) (375,0.456000) (376,0.454787) (377,0.453581) (378,0.452381) (379,0.453826) (380,0.452632) (381,0.454068) (382,0.452880) (383,0.454308) (384,0.455729) (385,0.457143) (386,0.455959) (387,0.454780) (388,0.456186) (389,0.457584) (390,0.456410) (391,0.457801) (392,0.456633) (393,0.458015) (394,0.456853) (395,0.458228) (396,0.457071) (397,0.458438) (398,0.459799) (399,0.458647) (400,0.457500) (401,0.456359) (402,0.455224) (403,0.456576) (404,0.455446) (405,0.456790) (406,0.455665) (407,0.457002) (408,0.455882) (409,0.454768) (410,0.453659) (411,0.452555) (412,0.453883) (413,0.455206) (414,0.456522) (415,0.457831) (416,0.459135) (417,0.458034) (418,0.459330) (419,0.458234) (420,0.459524) (421,0.458432) (422,0.457346) (423,0.458629) (424,0.457547) (425,0.456471) (426,0.457746) (427,0.456674) (428,0.457944) (429,0.456876) (430,0.455814) (431,0.454756) (432,0.456019) (433,0.454965) (434,0.456221) (435,0.455172) (436,0.456422) (437,0.457666) (438,0.458904) (439,0.457859) (440,0.459091) (441,0.460317) (442,0.459276) (443,0.460497) (444,0.459459) (445,0.458427) (446,0.457399) (447,0.458613) (448,0.459821) (449,0.458797) (450,0.460000) (451,0.458980) (452,0.460177) (453,0.459161) (454,0.460352) (455,0.459341) (456,0.460526) (457,0.461707) (458,0.460699) (459,0.459695) (460,0.460870) (461,0.462039) (462,0.461039) (463,0.462203) (464,0.461207) (465,0.462366) (466,0.463519) (467,0.462527) (468,0.463675) (469,0.462687) (470,0.461702) (471,0.462845) (472,0.463983) (473,0.465116) (474,0.466245) (475,0.467368) (476,0.466387) (477,0.465409) (478,0.466527) (479,0.465553) (480,0.464583) (481,0.463617) (482,0.462656) (483,0.461698) (484,0.460744) (485,0.459794) (486,0.460905) (487,0.462012) (488,0.463115) (489,0.462168) (490,0.461224) (491,0.462322) (492,0.461382) (493,0.462475) (494,0.463563) (495,0.462626) (496,0.463710) (497,0.462777) (498,0.463855) (499,0.462926) (500,0.464000) (501,0.465070) (502,0.466135) (503,0.467197) (504,0.466270) (505,0.467327) (506,0.468379) (507,0.469428) (508,0.470472) (509,0.471513) (510,0.470588) (511,0.471624) (512,0.470703) (513,0.469786) (514,0.470817) (515,0.471845) (516,0.472868) (517,0.473888) (518,0.474903) (519,0.473988) (520,0.475000) (521,0.476008) (522,0.475096) (523,0.474187) (524,0.473282) (525,0.472381) (526,0.473384) (527,0.474383) (528,0.475379) (529,0.476371) (530,0.475472) (531,0.474576) (532,0.473684) (533,0.472795) (534,0.471910) (535,0.472897) (536,0.473881) (537,0.472998) (538,0.473978) (539,0.473098) (540,0.472222) (541,0.473198) (542,0.474170) (543,0.475138) (544,0.474265) (545,0.475229) (546,0.474359) (547,0.475320) (548,0.476277) (549,0.475410) (550,0.476364) (551,0.477314) (552,0.478261) (553,0.479204) (554,0.480144) (555,0.479279) (556,0.478417) (557,0.477558) (558,0.476703) (559,0.477639) (560,0.478571) (561,0.479501) (562,0.478648) (563,0.477798) (564,0.478723) (565,0.479646) (566,0.480565) (567,0.481481) (568,0.480634) (569,0.479789) (570,0.480702) (571,0.481611) (572,0.480769) (573,0.479930) (574,0.479094) (575,0.480000) (576,0.480903) (577,0.480069) (578,0.480969) (579,0.481865) (580,0.482759) (581,0.481928) (582,0.481100) (583,0.480274) (584,0.479452) (585,0.480342) (586,0.479522) (587,0.480409) (588,0.479592) (589,0.478778) (590,0.479661) (591,0.478849) (592,0.479730) (593,0.478921) (594,0.478114) (595,0.478992) (596,0.478188) (597,0.477387) (598,0.476589) (599,0.477462) (600,0.476667) (601,0.475874) (602,0.476744) (603,0.475954) (604,0.476821) (605,0.476033) (606,0.475248) (607,0.476112) (608,0.476974) (609,0.476190) (610,0.477049) (611,0.477905) (612,0.478758) (613,0.479608) (614,0.480456) (615,0.481301) (616,0.480519) (617,0.481361) (618,0.480583) (619,0.481422) (620,0.482258) (621,0.483092) (622,0.482315) (623,0.483146) (624,0.482372) (625,0.483200) (626,0.484026) (627,0.483254) (628,0.484076) (629,0.484897) (630,0.485714) (631,0.484945) (632,0.484177) (633,0.483412) (634,0.484227) (635,0.483465) (636,0.484277) (637,0.485086) (638,0.484326) (639,0.485133) (640,0.485938) (641,0.485179) (642,0.485981) (643,0.486781) (644,0.486025) (645,0.486822) (646,0.487616) (647,0.486862) (648,0.487654) (649,0.488444) (650,0.487692) (651,0.486943) (652,0.486196) (653,0.485452) (654,0.486239) (655,0.485496) (656,0.484756) (657,0.485540) (658,0.484802) (659,0.484067) (660,0.484848) (661,0.485628) (662,0.484894) (663,0.484163) (664,0.484940) (665,0.485714) (666,0.484985) (667,0.485757) (668,0.485030) (669,0.484305) (670,0.485075) (671,0.484352) (672,0.483631) (673,0.482912) (674,0.482196) (675,0.482963) (676,0.483728) (677,0.484490) (678,0.483776) (679,0.483063) (680,0.482353) (681,0.483113) (682,0.483871) (683,0.483163) (684,0.482456) (685,0.481752) (686,0.482507) (687,0.481805) (688,0.481105) (689,0.480406) (690,0.479710) (691,0.480463) (692,0.481214) (693,0.480519) (694,0.481268) (695,0.482014) (696,0.481322) (697,0.480631) (698,0.479943) (699,0.480687) (700,0.481429) (701,0.482168) (702,0.481481) (703,0.482219) (704,0.481534) (705,0.480851) (706,0.481586) (707,0.480905) (708,0.480226) (709,0.480959) (710,0.481690) (711,0.482419) (712,0.481742) (713,0.482468) (714,0.483193) (715,0.483916) (716,0.483240) (717,0.483961) (718,0.483287) (719,0.482615) (720,0.483333) (721,0.482663) (722,0.483380) (723,0.482711) (724,0.482044) (725,0.481379) (726,0.482094) (727,0.482806) (728,0.483516) (729,0.482853) (730,0.482192) (731,0.482900) (732,0.482240) (733,0.481583) (734,0.482289) (735,0.482993) (736,0.482337) (737,0.483039) (738,0.482385) (739,0.481732) (740,0.481081) (741,0.481781) (742,0.482480) (743,0.481830) (744,0.482527) (745,0.483221) (746,0.483914) (747,0.484605) (748,0.485294) (749,0.484646) (750,0.485333) (751,0.486019) (752,0.485372) (753,0.486056) (754,0.485411) (755,0.484768) (756,0.484127) (757,0.484808) (758,0.485488) (759,0.486166) (760,0.486842) (761,0.487516) (762,0.488189) (763,0.488860) (764,0.489529) (765,0.488889) (766,0.489556) (767,0.490222) (768,0.490885) (769,0.491547) (770,0.490909) (771,0.490272) (772,0.489637) (773,0.490298) (774,0.490956) (775,0.490323) (776,0.490979) (777,0.490347) (778,0.489717) (779,0.489089) (780,0.488462) (781,0.487836) (782,0.488491) (783,0.487867) (784,0.487245) (785,0.486624) (786,0.486005) (787,0.485388) (788,0.484772) (789,0.484157) (790,0.483544) (791,0.482933) (792,0.482323) (793,0.481715) (794,0.482368) (795,0.483019) (796,0.482412) (797,0.483061) (798,0.483709) (799,0.483104) (800,0.483750) (801,0.483146) (802,0.483791) (803,0.483188) (804,0.482587) (805,0.481988) (806,0.481390) (807,0.480793) (808,0.481436) (809,0.482077) (810,0.482716) (811,0.482121) (812,0.482759) (813,0.483395) (814,0.484029) (815,0.484663) (816,0.485294) (817,0.484700) (818,0.485330) (819,0.485958) (820,0.485366) (821,0.484775) (822,0.485401) (823,0.486027) (824,0.485437) (825,0.486061) (826,0.486683) (827,0.487304) (828,0.486715) (829,0.487334) (830,0.486747) (831,0.486161) (832,0.485577) (833,0.484994) (834,0.485612) (835,0.485030) (836,0.484450) (837,0.485066) (838,0.485680) (839,0.486293) (840,0.486905) (841,0.487515) (842,0.486936) (843,0.487544) (844,0.488152) (845,0.487574) (846,0.486998) (847,0.487603) (848,0.487028) (849,0.486455) (850,0.485882) (851,0.486486) (852,0.487089) (853,0.487691) (854,0.488290) (855,0.487719) (856,0.487150) (857,0.486581) (858,0.487179) (859,0.486612) (860,0.486047) (861,0.486643) (862,0.486079) (863,0.485516) (864,0.484954) (865,0.485549) (866,0.486143) (867,0.485582) (868,0.485023) (869,0.484465) (870,0.483908) (871,0.484501) (872,0.485092) (873,0.484536) (874,0.485126) (875,0.485714) (876,0.486301) (877,0.486887) (878,0.487472) (879,0.488055) (880,0.487500) (881,0.488082) (882,0.487528) (883,0.488109) (884,0.487557) (885,0.488136) (886,0.488713) (887,0.489290) (888,0.489865) (889,0.489314) (890,0.488764) (891,0.488215) (892,0.488789) (893,0.488242) (894,0.488814) (895,0.489385) (896,0.489955) (897,0.490524) (898,0.489978) (899,0.490545) (900,0.491111) (901,0.490566) (902,0.490022) (903,0.489480) (904,0.490044) (905,0.489503) (906,0.488962) (907,0.489526) (908,0.490088) (909,0.490649) (910,0.490110) (911,0.489572) (912,0.490132) (913,0.489595) (914,0.490153) (915,0.489617) (916,0.489083) (917,0.488550) (918,0.489107) (919,0.488575) (920,0.489130) (921,0.488599) (922,0.489154) (923,0.488624) (924,0.489177) (925,0.489730) (926,0.490281) (927,0.489752) (928,0.490302) (929,0.490850) (930,0.491398) (931,0.491944) (932,0.492489) (933,0.493033) (934,0.492505) (935,0.493048) (936,0.493590) (937,0.494130) (938,0.493603) (939,0.493078) (940,0.492553) (941,0.492030) (942,0.491507) (943,0.492047) (944,0.491525) (945,0.491005) (946,0.490486) (947,0.489968) (948,0.490506) (949,0.491043) (950,0.491579) (951,0.492114) (952,0.492647) (953,0.492130) (954,0.491614) (955,0.491099) (956,0.490586) (957,0.490073) (958,0.490605) (959,0.490094) (960,0.489583) (961,0.490114) (962,0.490644) (963,0.491173) (964,0.490664) (965,0.491192) (966,0.490683) (967,0.491210) (968,0.491736) (969,0.492260) (970,0.492784) (971,0.493306) (972,0.493827) (973,0.494347) (974,0.493840) (975,0.493333) (976,0.492828) (977,0.492323) (978,0.492843) (979,0.493361) (980,0.492857) (981,0.493374) (982,0.493890) (983,0.494405) (984,0.494919) (985,0.494416) (986,0.494929) (987,0.495441) (988,0.494939) (989,0.494439) (990,0.493939) (991,0.494450) (992,0.493952) (993,0.494461) (994,0.494970) (995,0.495477) (996,0.495984) (997,0.495486) (998,0.495992) (999,0.495495) (1000,0.495000)};
\addplot[navy,dashed,domain=1:1000]{.5};\end{axis}\end{tikzpicture}\par\small Une trajectoire simulée, graine 2026 ; la droite est $y=\frac{1}{2}$.\end{center}
\item Notons $B_i$ l'indicatrice de pile. Alors $F_n=\frac{B_1+\cdots+B_n}{n}$. Chaque $B_i$ a espérance $\frac{1}{2}$ et variance $\frac{1}{4}$. Par linéarité, $E(F_n)=\frac{1}{2}$ ; par indépendance et changement d'échelle, $V(F_n)=\frac{n(\frac{1}{4})}{n^2}=\frac{1}{4n}$.
\item Bienaymé-Tchebychev donne $P(|F_{1000}-\frac{1}{2}|\ge0{,}05)\le\frac{1}{4\times1000\times0{,}05^2}=0{,}1$. Il s'agit d'une borne supérieure, pas de la probabilité exacte.
\item Pour assurer une probabilité complémentaire au plus $0{,}01$, il suffit de $\frac{1}{4n\times0{,}05^2}\le0{,}01$. Cela équivaut à $n\ge10000$. Alors $P(|F_n-\frac{1}{2}|<0{,}05)\ge0{,}99$.
\item La simulation illustre une réalisation finie ; d'autres graines donnent d'autres trajectoires. L'inégalité établit une garantie sur la probabilité, pour tout échantillon répondant aux hypothèses d'indépendance et de loi commune.
\end{enumerate}

\paragraph{E1.}La transformation affine donne $E(4X-1)=4E(X)-1=11$ et $V(4X-1)=4^2V(X)=32$. L'écart type est $\sqrt{32}=4\sqrt2$ ; le décalage $-1$ ne change pas la dispersion.

\paragraph{E2.}Par linéarité, $E(2X-Y)=2E(X)-E(Y)=4-5=-1$. L'indépendance de $X,Y$ permet d'additionner les variances de $2X$ et $-Y$ : $V(2X-Y)=4V(X)+V(Y)=4+4=8$.

\paragraph{E3.}La moyenne de $64$ variables indépendantes de même loi conserve l'espérance $10$. Sa variance vaut $\frac{4^2}{64}=\frac{1}{4}$, donc son écart type vaut $\frac{1}{2}$.

\paragraph{E4.}Bienaymé-Tchebychev avec $\delta=6$ donne $P(|X-100|\ge6)\le \frac{V(X)}{6^2}=\frac{9}{36}=\frac{1}{4}$. Cette borne ne donne pas une égalité ni une valeur exacte.

\paragraph{E5.}On a $P(|M_n-\mu|\ge0{,}2)\le\frac{4}{n\times0{,}2^2}=\frac{100}{n}$. Pour que cette borne soit au plus $0{,}05$, il suffit de $n\ge2000$. Le complément a alors une probabilité au moins $0{,}95$.

\subsection{Corrigé du DM14}\begin{enumerate}
\item Notons $B_i$ les indicatrices indépendantes des succès. $F_n=(\frac{1}{n})\sum_{i=1}^nB_i$, avec $E(B_i)=p$, $V(B_i)=p(1-p)$. Donc
\[E(F_n)=\frac1n\sum_{i=1}^n p=p,\qquad V(F_n)=\frac1{n^2}\sum_{i=1}^np(1-p)=\frac{p(1-p)}n.\]
\item $p(1-p)=p-p^2=\frac{1}{4}-(p-\frac{1}{2})^2\le\frac{1}{4}$. L'égalité est atteinte pour $p=\frac{1}{2}$.
\item Pour tout $p\in[0,1]$, $P(|F_n-p|\ge0{,}02)\le\frac{1}{4n\times0{,}02^2}$. Il suffit que cette borne soit au plus $0{,}05$, c'est-à-dire $n\ge12500$. Par complément, cette taille garantit la probabilité demandée, uniformément en $p$.
\item Sur $[0,0{,}1]$, la fonction $p\mapsto p(1-p)$ est croissante car sa dérivée $1-2p$ est positive. Donc $p(1-p)\le0{,}1\times0{,}9=0{,}09$. Il suffit désormais de $n\ge\frac{0{,}09}{0{,}05\times0{,}02^2}=4500$. Ces tailles sont suffisantes au sens de la majoration, pas nécessairement minimales pour la probabilité exacte.
\item Si tous les essais recopient une même variable $B$, alors $F_n=\frac{B+\cdots+B}{n}=B$. Donc $V(F_n)=p(1-p)$, sans division par $n$. L'addition des variances n'était permise que grâce à l'indépendance ; cette hypothèse disparaît ici. Pour $0<p<1$, augmenter le nombre de copies n'apporte aucune nouvelle information.
\end{enumerate}

\subsection{Corrigé du travail dans le manuel ou une annale}$M_N$ est le nombre moyen de donneurs universels par échantillon de $100$ personnes. La linéarité donne $E(M_N)=7{,}14$ ; l'indépendance donne $V(M_N)=\frac{6{,}63}N$.
\[P(|M_N-7{,}14|\ge0{,}14)\le\frac{6{,}63}{0{,}14^2N}.\]
Pour rendre cette borne au plus égale à $0{,}05$, il suffit de prendre
\[N\ge\left\lceil\frac{6{,}63}{0{,}05\times0{,}14^2}\right\rceil=6766.\]
Par complément, $P(7<M_N<7{,}28)\ge0{,}95$.

\fi

